\section{Introduction}
\label{sec:intro}

Dense 3D maps support navigation, manipulation, and scene understanding in robotics and augmented reality. TSDF methods~\cite{10.1145/3596711.3596726,newcombe2011kinectfusion} integrate range observations incrementally, while sparse volume structures~\cite{10.1145/2508363.2508374,Oleynikova_2017,vizzo2022vdbfusion} make large-scale onboard mapping practical. These systems accept poses from an external SLAM front-end. Loop closure and pose-graph optimization~\cite{Campos_2021,5979949} retroactively correct the trajectory, leaving the TSDF in the old world geometry and therefore inconsistent with the revised poses (Fig.~\ref{fig:teaser}).

% \vspace{1mm}\noindent\textit{The map-rebuild bottleneck.}
The standard remedy is to discard the map and reintegrate all $N$ range frames with corrected poses~\cite{Cramariuc_2023,rosinol2021kimeraslamspatialperception}. Its $O(N|\Omega|)$ cost, for $|\Omega|$ measurements per frame, is paid again after every global correction and can reach minutes on long sessions. A de-integration/reintegration cycle~\cite{han2018flashfusion} avoids a full rebuild but integrates twice for each affected frame; after a global correction nearly every frame is affected, so its cost upper-bounds full reintegration rather than reducing it. We therefore target a narrow interface: consume an already fused TSDF and a revised trajectory, and return a TSDF in the same format that the mapper keeps fusing into.
 

Submap methods such as Voxgraph~\cite{reijgwart2020voxgraph} provide this interface but require the map to be partitioned during fusion, and their rigid per-submap transforms leave intra-submap drift unresolved. The remaining approaches change the representation instead: non-rigid surfel SLAM~\cite{whelan2016elasticfusion} corrects its own tracking state rather than an external TSDF, and neural implicit maps~\cite{sucar2021imap,zhu2022niceslam,liso2024loopyslam,bruns2024neuralgraphmap} require representation-specific re-optimization. Each is a sound design; none returns a corrected TSDF to a mapper already running one.

Our starting point is an invariance: a calibrated range sample expressed in its sensor frame is unchanged by a later world-frame pose correction, although its relation to the current map surface is not. WarpVox therefore uses keyframe range measurements to rebuild projective measurement--surface correspondences at each correction event. Re-expressing the matched samples with corrected poses yields corrected world-space surface targets by projection alone: the correction cost then follows the current map rather than the number of measurements behind it or the volumetric work of fusing them.

WarpVox combines four stages: (i) projectively re-associate keyframe range measurements with the current pre-correction surface, (ii)~compute robust targets from the corrected poses, (iii)~propagate sparse displacements through a single ARAP solve~\cite{ARAP_modeling:2007} on a QEM proxy, and (iv)~reconstruct a TSDF whose zero set follows the warped surface while observed support, sign, and weight primarily follow the source volume. At each loop-closure, this corrected volume replaces the inconsistent active TSDF; Voxblox then resumes fusion from the updated state until a later trajectory revision triggers the same correction procedure again. Only Stages~1--2 use keyframe measurements to construct correction targets; deformation and reconstruction operate on the active geometry and output volume.

\begin{itemize}
% \item We introduce an event-level system that converts an external trajectory revision and an already fused active TSDF into a corrected, fusion-ready TSDF without replaying the complete historical range stream.
% \item The system constructs correction targets through projective measurement--surface reassociation, propagates them through one proxy deformation, and reconstructs source-supported TSDF fields so that the corrected volume can replace the active map.
\item An observation-derived target construction that turns cached keyframe range measurements and a revised trajectory into corrected surface targets by projection, without replaying volumetric integration. 
\item A source-supported TSDF reconstruction that pairs the warped zero set with support, sign, and integration weights pulled back from the source volume, returning a native Voxblox map that replaces the active TSDF and continues fusing.
\item An evaluation of RGB-D and LiDAR correction against reintegration and mapping baselines over 51 loop-closure cases, together with a separate comparison of final-state surface correction against Kimera-PGMO on ten complete bags.
\end{itemize}
